\documentclass[11pt]{article}
\usepackage[margin=1in]{geometry}
\usepackage{amsmath,amssymb,amsthm}
\usepackage[hidelinks]{hyperref}

\newtheorem*{theorem}{Literature Answer}

\title{Positive Literature Answer to the Canonical Orbit-Pseudometric Question in arXiv:1804.11164}
\author{FA/Banach run packet}
\date{2026-07-18}

\begin{document}
\maketitle

\section*{Status}

This is a literature-already-answered packet. It records that a question from
the concluding problems of Cuth--Doucha--Kurka \cite{CDK1} was answered
positively in the later paper of Kurka \cite{Kurka}. It is not an original proof
from this run.

\section*{Original Problem}

In the concluding remarks of \cite{CDK1}, the authors ask for a continuous
version of the Clemens--Gao--Kechris theorem. Let \(d\) be the Hausdorff
distance on \(F(\mathbb U)\), let \(\operatorname{Iso}(\mathbb U)\) act
canonically on \(F(\mathbb U)\), and let \(\rho\) be the corresponding orbit
pseudometric. The question is whether
\[
  \rho_{GH}
  \quad\text{and}\quad
  \rho
\]
are Borel-uniformly continuously bi-reducible, or whether at least one of the
two reductions holds.

\section*{Supporting Answer}

Kurka \cite[Theorem 1.2]{Kurka} proves precisely that
\[
  \rho_{GH}
  \equiv_{B,u}
  \rho_{\operatorname{Iso}(\mathbb U),\rho_H}
\]
on \(F(\mathbb U)\setminus\{\emptyset\}\), where \(\rho_H\) is the Hausdorff
distance. This is the same canonical Urysohn-space orbit pseudometric asked
about in \cite{CDK1}, restricted to the usual nonempty closed-set coding of
Polish metric spaces.

The same paper proves a stronger comparison theorem. In Theorem 1.3 it
constructs a Borel map \(p\mapsto Y_p\) such that
\[
  \rho_{GH}(Y_p,Y_q)
  \leq \rho_{\operatorname{Iso}(\mathbb U),\rho_H}(Y_p,Y_q)
  \leq \rho_{G,d}(p,q)
  \leq 2\rho_{GH}(Y_p,Y_q).
\]
Consequently the reductions in Theorem 1.2 can be chosen Borel-Lipschitz on
small distances.

\begin{theorem}
The canonical orbit-pseudometric question in \cite{CDK1} has a positive
literature answer: the Gromov--Hausdorff pseudometric is Borel-uniformly
continuously bi-reducible with the orbit pseudometric induced by the canonical
action of \(\operatorname{Iso}(\mathbb U)\) on \(F(\mathbb U)\) equipped with
Hausdorff distance.
\end{theorem}

\begin{proof}[Literature identification]
The question in \cite{CDK1} uses exactly the data
\((F(\mathbb U),\operatorname{Iso}(\mathbb U),\rho_H)\). Kurka's Theorem 1.2
states the Borel-u.c. bireducibility of \(\rho_{GH}\) and
\(\rho_{\operatorname{Iso}(\mathbb U),\rho_H}\) for this data. Therefore both
directions requested by the question are supplied by the later paper.
\end{proof}

\section*{Scope Limitations}

This packet answers only the specific canonical Urysohn/Hausdorff
orbit-pseudometric question. It does not solve the other concluding problems
from \cite{CDK1}; in particular, Kurka leaves open a broader question about
removing the auxiliary system assumption in the general orbit-pseudometric
reduction theorem.

\section*{Verification Notes}

The source and supporting papers are distinct. The supporting paper cites
\cite{CDK1} and explicitly presents Theorem 1.2 as a positive answer to
Question 41 from that paper. The local duplicate search covered arXiv ids
\texttt{1804.11164} and \texttt{2204.08375}, plus phrases
\texttt{Gromov-Hausdorff}, \texttt{orbit pseudometric},
\texttt{Iso(U)}, \texttt{Hausdorff distance}, and \texttt{bireducible}.

\begin{thebibliography}{9}

\bibitem{CDK1}
M. Cuth, M. Doucha, and O. Kurka,
\emph{Complexity of distances: Theory of generalized analytic equivalence relations},
arXiv:1804.11164; J. Math. Log. 22 (2022), no. 1, 2250005.

\bibitem{Kurka}
O. Kurka,
\emph{Orbit pseudometrics and a universality property of the Gromov-Hausdorff distance},
arXiv:2204.08375; Topology Appl. 364 (2025), 109095.

\end{thebibliography}

\end{document}
